\scanpage{56}
\mspage{43}
% archive sheet number: 55

y aurait lieu d'introduire la notion de $C$-carré

\diagram{\begin{tikzcd}[ampersand replacement=\&,column sep=small,row sep=tiny]
 X_0\arrow[r,dash]\arrow[d,dash] \& X_1\arrow[d,dash]\\
 X_{\unsure{2|\dot{2}}}\arrow[r,dash] \& \unsure{x|X}_3
\end{tikzcd}}

, comme un carré tel que $X_{\unsure{1}}\amalg_{X_0}X_2\unsure{\mathbin{-}|\longrightarrow}X_3$ soit dans $C$\inserted{(NB notion symétrique pour les \unsure{opérations|foncteurs} de symétrie des carrés)}, et dire que le composé de deux $C$-carrés est un $C$-carré. $]$

Cochangement de base

\diagram{\begin{tikzcd}[ampersand replacement=\&,row sep=small]
 X_{\unsure{\cdot|\bullet}}\arrow[r,"u"]\arrow[d] \& Y_{\cdot}\arrow[d,dash]\\
 X'_{\cdot}\arrow[r,dash,"u'"] \& Y'_{\cdot}
\end{tikzcd}}

, Considérons $\deleted{\unsure{J}}K\subset J$, on a

\diagram{\begin{tikzcd}[ampersand replacement=\&,column sep=small,row sep=small]
 \& X_K\arrow[r]\arrow[d,dash]\arrow[ddl,thick] \& X_{\deleted{\unsure{K}}J}\arrow[d]\arrow[dr,thick]\arrow[ddl,thick] \&\\
 \& Y_K\arrow[r]\arrow[ddl]\arrow[ur,phantom,"{\text{\unsure{cocart}}}" description] \& Y_{K,J}\arrow[r,dash,"{u_{K,J}}"']\arrow[ddl] \& Y_J\arrow[ddl,thick]\\
 X'_K\arrow[r,leftrightarrow,thick]\arrow[d,thick]\arrow[ur,phantom,"{\text{cocart}}" description,sloped] \& X'_J\arrow[d,thick]\arrow[dr,thick] \& \&\\
 Y'_K\arrow[r,leftrightarrow,thick]\arrow[ur,phantom,"{\text{cocart}}" description] \& Y'_{J,K}\arrow[r,thick,"{u'_{\unsure{K},J}}"] \& Y'_J\arrow[uur,phantom,"{\text{cocart}}" description,sloped] \&
\end{tikzcd}}

, où les trois carrés \inserted{\unsure{latéraux}} marqués ``cocart'' sont cocartésiens, ainsi que le carré en plan incliné. Il en résulte, d'abord que le quatrième carré latéral du cube est lui aussi cocartésien,
